% !TEX root = ../main.tex

\section{Panoramica sugli UAV VTOL}

\subsection{Configurazioni principali}
La categoria degli UAV a decollo e atterraggio verticale (VTOL) è vasta e si articola principalmente in tre configurazioni strutturali, ognuna con specifici vantaggi e limitazioni operative:
\begin{itemize}
    \item \textbf{Multirotori:} Questa classe è caratterizzata da una notevole semplicità meccanica e dalla capacità di mantenere un \textit{hovering} estremamente stabile. Tuttavia, soffrono di una bassa efficienza energetica nel volo traslato, cosa che ne limita notevolmente l'autonomia e il raggio d'azione.
    \item \textbf{Ala fissa:} I velivoli ad ala fissa convenzionali offrono un'elevata efficienza aerodinamica, permettendo lunghe percorrenze e velocità di crociera elevate. Il loro svantaggio principale risiede nella necessità di infrastrutture dedicate o spazi ampi per le manovre di decollo e atterraggio, precludendone l'uso in ambienti confinati.
    \item \textbf{Ibridi:} Questa categoria nasce per combinare la versatilità operativa dei multirotori con l'efficienza dei velivoli ad ala fissa. Include architetture complesse come i \textit{tilting-rotor}, \textit{tilting-wing} e \textit{tailsitter}, capaci di transire tra le modalità di volo verticale e orizzontale.
\end{itemize}

\subsection{Sfide di modellazione e controllo}
Lo sviluppo di velivoli VTOL ibridi presenta sfide ingegneristiche significative, concentrate prevalentemente nella gestione della fase di transizione tra il volo verticale e quello orizzontale. In questo regime, il velivolo è soggetto a forti non linearità aerodinamiche causate dall'interazione complessa tra i flussi dei rotori, le superfici alari e la fusoliera.
Una modellazione efficace deve catturare accuratamente fenomeni quali gli effetti giroscopici, le variazioni repentine di portanza e resistenza, e il forte accoppiamento tra i gradi di libertà longitudinale e laterale. 
Di conseguenza, i sistemi di controllo devono essere progettati per garantire robustezza e prestazioni elevate durante tutte le fasi di volo, adattandosi a condizioni operative che variano drasticamente tra l'hovering statico e la crociera aerodinamica.

\clearpage
\section{UAV VTOL con ali indipendentemente inclinabili}
Un contributo significativo nello sviluppo di configurazioni ibride innovative è stato presentato in \cite{nam2020}, dove viene proposto un UAV VTOL caratterizzato da ali inclinabili in modo indipendente. L'obiettivo primario di tale design è superare i limiti di manovrabilità dei tricotteri classici e mitigare la coppia di reazione tipicamente gestita dal rotore di coda.
\\Dal punto di vista della modellazione, è stato derivato un modello dinamico a sei gradi di libertà (6-DOF) basato sulla formulazione di Newton-Eulero. Il modello evidenzia come il controllo differenziale dell'angolo di inclinazione (\textit{tilt}) delle due ali generi una differenza di portanza in grado di indurre momenti di rollio, sfruttando le forze propulsive in combinazione con quelle aerodinamiche.
\\La strategia di controllo implementata si basa su regolatori PID e PI integrati nel firmware Ardupilot, strutturati in base alla fase di volo:
\begin{itemize}
    \item \textbf{Modalità Verticale:} Vengono impiegati PID classici per la stabilizzazione di rollio, beccheggio e imbardata, allocando i comandi su tre motori brushless e tre servocomandi.
    \item \textbf{Modalità Orizzontale:} Il controllo del rollio è affidato alla variazione differenziale del tilt alare, il beccheggio è gestito tramite il rotore di coda (spinta e inclinazione vettoriale), mentre l'imbardata è stabilizzata da un controllore PI che minimizza lo scivolamento laterale (\textit{sideslip}).
\end{itemize}
La transizione è governata da una logica a stati finiti che coordina l'inclinazione progressiva delle ali, garantendo una traiettoria fluida durante il cambio di configurazione.

\clearpage
\section{Modellazione e controllo Backstepping per Quadricotteri}
In \cite{mian2008}, viene affrontato il problema del controllo non lineare per quadricotteri attraverso la derivazione di un modello completo a 6-DOF che include effetti giroscopici e dinamica degli attuatori. L'analisi sottolinea la natura sottattuata e fortemente accoppiata del sistema.
La strategia di controllo proposta adotta un approccio sequenziale ibrido:
\begin{itemize}
    \item Per il \textbf{sottosistema traslazionale} (posizione e altitudine), viene utilizzata una tecnica di linearizzazione con retroazione accoppiata a un regolatore PD.
    \item Per il \textbf{sottosistema rotazionale} (assetto), viene implementato un controllore non lineare basato sulla tecnica \textit{Backstepping} integrata con un'azione integrale (PID).
\end{itemize}
L'uso del Backstepping garantisce la stabilità asintotica globale, dimostrata tramite funzioni di Lyapunov, e offre una robustezza superiore rispetto ai PID tradizionali, specialmente nella reiezione dei disturbi esterni e nella riduzione dell'errore di inseguimento a regime.

\clearpage
\section{Controllo Sliding Mode "Smoothed" e Non-Overshooting}
Un'evoluzione recente per affrontare simultaneamente il problema del \textit{chattering} e quello delle sovraelongazioni (\textit{overshoot}) nella risposta dinamica degli UAV è presentata in \cite{Wang2024}.
Gli autori propongono una strategia di controllo robusta denominata \textit{Smoothed Non-overshooting Sliding Mode Control}, progettata per garantire la stabilità esponenziale globale anche in presenza di disturbi stocastici limitati.

A differenza degli approcci classici, che utilizzano una funzione di saturazione lineare a tratti ($\text{sat}(s)$) per mitigare il chattering, questo metodo impiega una funzione basata sulla tangente iperbolica ($\tanh(\cdot)$).
La legge di controllo assume la forma generale:
\begin{equation}
    u(t) = -k \cdot \tanh(\rho \cdot s(t))
\end{equation}
dove $\rho$ è un parametro di sintonizzazione che determina la pendenza della funzione.
I vantaggi principali di questa formulazione includono:
\begin{itemize}
    \item \textbf{Continuità e Derivabilità:} L'uso della $\tanh(\cdot)$ rende il segnale di controllo una funzione liscia (\textit{smooth}) e ovunque differenziabile. Questo elimina le discontinuità nella derivata che si verificano ai bordi dello strato limite nella funzione $\text{sat}(\cdot)$, riducendo lo stress meccanico sugli attuatori e migliorando la precisione di tracking.
    \item \textbf{Approssimazione Universale:} Viene dimostrato che per valori sufficientemente elevati di $\rho$, la funzione $\tanh(\rho \cdot s)$ approssima la funzione $\text{sign}(s)$ ideale, permettendo di recuperare le proprietà di robustezza dello SMC puro pur mantenendo un comportamento continuo.
    \item \textbf{Assenza di Overshoot:} La struttura del controllore è progettata con due sottosistemi sequenziali (raggiungibilità e stabilità) che garantiscono che la variabile di errore converga a zero senza mai superare il target, una caratteristica critica per la sicurezza nelle manovre UAV in ambienti ostili.
\end{itemize}

\clearpage
\section{Modellazione e Controllo Robusto per UAV Coaxial Tilting-Rotor}
Un'evoluzione significativa nel panorama dei velivoli a decollo verticale è rappresentata dalla configurazione \textit{Coaxial Tilting-Rotor} (CTRUAV), analizzata in \cite{cheng2022}. Questa architettura si distingue per la presenza di due coppie di rotori coassiali inclinabili posizionati nella parte anteriore del velivolo e un rotore fisso in coda. Rispetto ai tricotteri convenzionali, l'impiego di rotori coassiali permette di generare una spinta maggiore senza aumentare le dimensioni del telaio, migliorando la manovrabilità e annullando naturalmente il momento giroscopico grazie alla configurazione controrotante.
\begin{figure}[h]
    \centering
    \includegraphics[width=0.5\textwidth]{coaxial}
    \caption{Coaxial tilting rotor UAV.}
    \label{fig:coaxial}
\end{figure}
Il modello dinamico a sei gradi di libertà è derivato utilizzando le formule di Newton-Eulero. Per governare questo sistema in presenza di disturbi esterni, è stata proposta una strategia di controllo gerarchica:
\begin{itemize}
    \item \textbf{Loop Esterno (Posizione):} Per il tracking della traiettoria desiderata, viene impiegato un controllore basato sulla tecnica del \textit{Backstepping}, che stabilizza asintoticamente l'errore di posizione.
    \item \textbf{Loop Interno (Assetto):} La stabilizzazione è affidata a un controllore \textit{Robust Adaptive Backstepping Sliding Mode} (ABSMC). L'integrazione di una legge adattiva permette di stimare online il limite superiore dei disturbi esterni. Tale stima modula il guadagno della legge di controllo, minimizzando il fenomeno del \textit{chattering} tramite l'uso di funzioni di saturazione.
\end{itemize}
Le simulazioni dimostrano che l'approccio proposto garantisce una robustezza superiore rispetto ai controllori backstepping tradizionali.

\clearpage
\section{Strategie di Volo e Controllo per Tricotteri Tilting-Rotor}
In \cite{chen2017}, l'attenzione si sposta sulla realizzazione pratica di un volo autonomo "full-mode" per un UAV tilting-rotor in configurazione tricottero compatta. L'obiettivo è superare le limitazioni di costo associate all'ottenimento di modelli dinamici precisi, proponendo una strategia robusta che non dipenda da dati aerodinamici ad alta fedeltà.
\\La gestione del volo è affidata a una strategia basata su una macchina a stati finiti che coordina le diverse fasi: \textit{Rotor mode} (decollo verticale), \textit{Conversion} (transizione accelerata), \textit{Fixed-wing mode} (crociera) e \textit{Reconversion} (ritorno all'hovering).
Un elemento centrale è il sistema di \textbf{Motor Mixing Algorithm}; poiché il sistema è sovra-attuato e non affine durante la transizione, il mixer distribuisce i pesi di controllo tra modalità rotore e ala fissa in base all'efficienza degli attuatori e all'angolo di tilting corrente. Una peculiarità è il controllo dell'imbardata in hovering tramite tilt differenziale dei rotori anteriori, che viene disattivato progressivamente durante la transizione.
I test di volo reali hanno validato la capacità del velivolo di eseguire l'intero inviluppo di volo mantenendo la stabilità \cite{chen2017}.
\begin{figure}[h]
        \centering
        \includegraphics[width=0.5\textwidth]{paper5.png}
        \caption{Control logic in rotor mode. (a) roll, (b) pitch, (c) thrust, and (d) yaw.}
        \label{fig:paper5}
\end{figure}

\clearpage
\section{Controllo Sliding Mode a Tempo Fisso}
La problematica della convergenza rapida e garantita è affrontata in \cite{yu2022}, dove viene proposto uno schema per la stabilizzazione dell'assetto di un \textit{Tilting Trirotor UAV} basato sulla stabilità a tempo fisso (\textit{Fixed-Time Stability}). A differenza della stabilità asintotica o a tempo finito, questo approccio garantisce che il sistema raggiunga l'equilibrio entro un tempo limitato, noto a priori e indipendente dallo stato iniziale.
\\Il contributo teorico principale è la formulazione di un \textit{Novel Fixed-Time Stable System} che dimostra un tasso di convergenza superiore rispetto ai metodi esistenti. Basandosi su questo, è stato sviluppato uno schema \textit{Nonsingular Fixed-Time Sliding Mode Control} (NFTSMC) che include:
\begin{itemize}
    \item \textbf{Superficie di Scorrimento Non Singolare:} Progettata per evitare le singolarità matematiche tipiche dei controllori terminali e garantire la convergenza.
    \item \textbf{Tecnica Adattiva per i Disturbi:} Una legge di aggiornamento stima i disturbi in tempo reale, assicurando robustezza.
\end{itemize}
La validazione sperimentale su prototipo reale ha confermato una convergenza degli errori di assetto in meno di 0,3 secondi e una notevole reiezione dei disturbi \cite{yu2022}.

\clearpage
\section{Controllo tramite Spinta Differenziale}
Un aspetto tecnologico cruciale esplorato in \cite{riboldi2023} è la completa sostituzione delle superfici di controllo aerodinamiche convenzionali (alettoni, elevatori, timoni) con una logica di controllo basata esclusivamente sulla modulazione della spinta (\textit{Differential Thrust}).
Questa architettura richiede l'impiego di propulsori in grado di fornire spinta reversibile (positiva e negativa), permettendo la generazione di momenti di controllo anche a velocità nulla, una capacità preclusa ai sistemi aerodinamici tradizionali.
\\La strategia di (\textit{mixing}) mappa i comandi di assetto desiderati (Rollio, Beccheggio, Imbardata) direttamente sulle coppie di motori, secondo lo schema seguente (riferito a una configurazione a quattro propulsori disposti a "X"):
\begin{itemize}
    \item \textbf{Controllo di Beccheggio (Pitch):} Sostituisce la funzione dell'elevatore. Viene generato creando una spinta differenziale tra i rotori superiori e quelli inferiori. Ad esempio, per cabrare, i propulsori inferiori generano spinta positiva mentre quelli superiori generano spinta negativa, creando una coppia pura sull'asse trasversale.
    \item \textbf{Controllo di Imbardata (Yaw):} Sostituisce il timone verticale. Si ottiene differenziando la spinta tra i motori di destra e quelli di sinistra. Una rotazione verso destra è indotta da una spinta positiva sui motori di sinistra e negativa su quelli di destra.
    \item \textbf{Controllo di Rollio (Roll):} Sostituisce gli alettoni. A differenza dei velivoli ad ala fissa dove il rollio è generato alle estremità alari, qui viene prodotto tramite una configurazione incrociata. Ad esempio, una spinta positiva del motore in basso a destra combinata con il motore in alto a sinistra, opposta alla coppia degli altri due, genera il momento torcente longitudinale necessario \cite{riboldi2023}.
\end{itemize}
Questa metodologia offre il vantaggio di eliminare attuatori meccanici complessi (servocomandi) e parti mobili, riducendo il peso strutturale e la manutenzione. Tuttavia, introduce una maggiore complessità nel software di controllo, poiché i gradi di libertà risultano fortemente accoppiati: una singola manovra richiede spesso la modulazione simultanea di tutti i propulsori per evitare effetti secondari indesiderati sugli altri assi.